% introduction.tex -- Introduction

Generating novel protein sequences that respect the statistical constraints of a family is a central problem in computational biology, with applications in enzyme engineering, therapeutic design, and understanding molecular evolution~\cite{wuMachineLearningAidedProtein2019}. Modern generative approaches---variational autoencoders~\cite{kingmaWellingAutoEncoding2014,sinaiVariationalAutoEncoding2017}, autoregressive language models~\cite{madaniProGenLanguageModeling2023,ferruzProtGPT2DeepUnsupervised2022}, diffusion models~\cite{alamdariProteinGeneration2023}, and MSA-conditioned transformers~\cite{raoMSATransformer2021}---have demonstrated impressive performance when trained on thousands to millions of sequences. However, most protein families are small: the median Pfam family contains fewer than 100 full-length members~\cite{mistryPfamProteinFamilies2021}. In this scarce-data regime, learned generative models overfit, collapse in diversity, or fail to train altogether.

We propose a fundamentally different approach. Rather than learning a generative model, we observe that the modern Hopfield energy~\cite{ramsauerHopfieldNetworksAll2021} over a set of stored protein sequences defines a Boltzmann distribution whose score function---the gradient of the log-density---is given in closed form by the residual of a single attention operation. Langevin dynamics on this energy yields a training-free sampler that we call \emph{stochastic attention}: each iteration performs a softmax-weighted pull toward stored family members, a contraction, and an isotropic Gaussian perturbation, requiring only two matrix-vector products per step.

The key question this paper addresses is: \emph{what can you generate from 50--200 protein sequences with no training?} We show that stochastic attention, applied to aligned sequences from multiple Pfam families and antimicrobial peptide datasets, produces novel sequences that preserve family-level amino acid composition, positional conservation, and pairwise couplings, while achieving high structural plausibility as assessed by AlphaFold2~\cite{jumperHighlyAccurateProtein2021}. We characterize the critical inverse temperature $\beta^*$ at which the sampler transitions from diffuse exploration to pattern retrieval, and show that $\beta^*$ can be predicted directly from multiple sequence alignment (MSA) statistics without grid search. A systematic scaling study reveals that stochastic attention maintains generation quality down to $K=50$ family members, a regime where protein-specific learned baselines---DeepSequence~\cite{riesselman2018deep}, EvoDiff~\cite{alamdariProteinGeneration2023}, ProtGPT2~\cite{ferruzProtGPT2DeepUnsupervised2022}, and MSA Transformer~\cite{raoMSATransformer2021}---degrade substantially.
